\documentclass{beamer}
	
	\usepackage[utf8]{inputenc} 
	\usepackage[T1]{fontenc}      
	\usepackage[francais]{babel}
	\usepackage[absolute,overlay]{textpos}
	\usepackage{graphicx}
	
	\graphicspath{figures/}
	
	\useoutertheme{shadow}% Pied de page
	\usetheme{PaloAlto}% Thème avec sidebar
	
	%\setbeamercolor{structure}{fg=cyan}
	
	% Mise en forme pour la première diapo
	\title[Présentation]{Retrogaming : Metal Slug}
	\subtitle[\ldots]{Soutenance Projet}
	\author[BONNIN GIRARD SILOTIA WAGNER]{Thomas BONNIN Jules GIRARD \\ Gilles SILOTIA Pascal WAGNER}
	\institute[UM2]{Université de Montpellier}
	\institute[UM2]{Université de Montpellier \\ Faculté des Sciences}
	\date{\today}
	
	% Retire les noms et titre de la sidebar
	\makeatletter
		\setbeamertemplate{sidebar \beamer@sidebarside}%{sidebar theme}
		{
			\beamer@tempdim=\beamer@sidebarwidth%
			\advance\beamer@tempdim by -6pt%
			\insertverticalnavigation{\beamer@sidebarwidth}%
			\vfill
			\ifx\beamer@sidebarside\beamer@lefttext%
			\else%
				\usebeamercolor{normal text}%
				\llap{\usebeamertemplate***{navigation symbols}\hskip0.1cm}%
				\vskip2pt%
			\fi%
		}%
	\makeatother
	
	% Numéro de diapo (en com rajoute le nombre de diapos au total -> 1/20 ... 20/20)
	\addtobeamertemplate{footline}{\insertframenumber}%/\inserttotalframenumber}
	
	% Insère une diapo à chaque changement de section comportant le sommaire
	% avec la nouvelle section courante (et ses sous-sections) en noir et
	% le reste en grisé/transparent
	\AtBeginSection[]
	{
		\begin{frame}
			\frametitle{Sommaire}
			\tableofcontents[currentsection, hideothersubsections]
		\end{frame} 
	}

	\begin{document}
		
		% Première diapo sans sidebar [plain] et centrer
		\begingroup
			\makeatletter
				\setlength{\hoffset}{-.5\beamer@sidebarwidth}
			\makeatother
			\begin{frame}[plain]
				\begin{textblock*}{2cm}(2cm,6cm) % {block width} (coords)
					\includegraphics[width=2cm]{figures/logo_fds.png}
				\end{textblock*}
				\begin{textblock*}{2.25cm}(10cm,6cm) % {block width} (coords)
					\includegraphics[width=2.25cm]{figures/logo_um.png}
				\end{textblock*}
				\titlepage
			\end{frame}
		\endgroup
		
		% Deuxième diapo annonce du plan/sommaire
		\begin{frame}
			% Titre diapo
			\frametitle{Sommaire}
			% Plan sans sous-sections
			\tableofcontents[hideallsubsections]
		\end{frame}
		
		
		%%%%%%%%%%%%%%%%%%%%%%%%%%%%%%%%%%%%%%%%%%%%%%%%%%%%%%%%%%%%%%
		%%%														   %%%
		%%%	Les diapos qui suivent (avant la première section) ne  %%%
		%%%	sont que des tests et des exemples pour nous aider     %%%
		%%%														   %%%
		%%%%%%%%%%%%%%%%%%%%%%%%%%%%%%%%%%%%%%%%%%%%%%%%%%%%%%%%%%%%%%
		
		
		% Quelques mises en forme avec des beaux cadres
		\begin{frame}
			\begin{block}{Qu'est-ce un bloc sous Beamer ?}
				Tout simplement ceci !
			\end{block}
			\begin{definition} % Définition
				On donne une définition.
			\end{definition}

			\begin{example} % Exemple
				Suivi d'un exemple.
			\end{example}

			\begin{theorem} % Théorème
				On peut éventuellement ajouter un théorème.
			\end{theorem}

			\begin{proof} % Démonstration
				Suivi de sa démonstration.
			\end{proof}
		\end{frame}
		
		
		\begin{frame}
			\frametitle{Un exemple de titre}
			\framesubtitle{avec un exemple de sous-titre}
			Enfin, le texte ! :)
		\end{frame}
		
		% Considérer comme une même diapo mais avec des animations
		\begin{frame}
			\begin{itemize}
				\item<1,3> Ploum!
				\item<2-4> Plim?
				\item Plum...
			\end{itemize}
		\end{frame}
		
		
		
		
		
		
		% Introduction
		\section{Section 1}
		
		% Cahier des charges
		\section{Section 2}
		
		% Analyse
		\section{Section 3}
		
		% Conception/réalisations
		\section{Section 4}
		
		% UNE AUTRE ??? 
		% Améliorations/Perspectives
		\section{Section 5}
		
		% Conclusion
		\section{Section 6}
		
		
		
		
	\end{document}
